\section*{Chapitre \ref{ch:b1:cell-unit-of-life} --- La cellule~: unité du vivant}

\begin{solution}{exo:b1:cell-unit-of-life:1}
Tous les \omterm{def:b1:organism-environment:organism}{organismes} sont faits de \omterm{def:b1:cell-unit-of-life:cell}{cellules} (tous les \omterm{def:b1:body-plans-tissues:tissue}{tissus} examinés au
microscope, végétaux comme animaux, en montrent)~; toute \omterm{def:b1:cell-unit-of-life:cell}{cellule} provient
d'une \omterm{def:b1:cell-unit-of-life:cell}{cellule} (division observée dans les embryons et chez les algues~;
les ballons de Pasteur excluent la génération spontanée)~; le matériel
héréditaire passe à la division (chromosomes suivis au cours de la
mitose).
\end{solution}

\begin{solution}{exo:b1:cell-unit-of-life:2}
\omterm{def:b1:cell-unit-of-life:prokeuk}{Noyau} (absent / présent), \omterm{def:b1:cell-unit-of-life:prokeuk}{organites} limités par une membrane (absents /
présents), chromosome (un seul, circulaire / plusieurs, linéaires),
taille (\qtyrange{0.5}{5}{\micro m} / \qtyrange{10}{100}{\micro m}),
\omterm{def:b1:cell-unit-of-life:prokeuk}{paroi} de peptidoglycane (bactéries) / de cellulose, de chitine ou
absente, cytosquelette (rudimentaire / élaboré).
\end{solution}

\begin{solution}{exo:b1:cell-unit-of-life:3}
$\qty{1}{mm}/\qty{1}{\micro m} = 1000$ bactéries~; $1000/20 = 50$
\omterm{def:b1:cell-unit-of-life:cell}{cellules} animales.
\end{solution}

\begin{solution}{exo:b1:cell-unit-of-life:4}
La résolution dépend de la longueur d'onde et de l'ouverture numérique,
non du \omterm{def:b1:cell-unit-of-life:resolution}{grossissement}~; au-delà de $1000\times$ environ, l'image agrandit
le flou de la limite de diffraction (\qty{0.2}{\micro m}) sans rien
séparer de nouveau.
\end{solution}

\begin{solution}{exo:b1:cell-unit-of-life:5}
$0.61\times 550/0.65 = \qty{516}{nm}$~; dans le bleu, \qty{422}{nm}.
Deux bactéries distantes de \qty{500}{nm}~: tout juste, en lumière
bleue~; des ribosomes distants de \qty{30}{nm}~: non, d'un facteur
quinze.
\end{solution}

\begin{solution}{exo:b1:cell-unit-of-life:6}
Violettes~: \omterm{def:b1:cell-unit-of-life:prokeuk}{paroi} épaisse de peptidoglycane retenant le complexe violet
de cristal--iode, sans membrane externe (Gram positives). Roses~:
peptidoglycane mince sous une membrane externe, colorant éliminé au
rinçage, contre-coloration (Gram négatives). \emph{E. coli} est un
bâtonnet Gram négatif~: les bâtonnets roses.
\end{solution}

\begin{solution}{exo:b1:cell-unit-of-life:7}
Mitochondrie~: $v = \frac{2}{9}\times(5\times 10^{-7})^2\times 100\times
9.81\times 10^4/10^{-3} = \qty{5.4e-4}{m/s}$~; \qty{5}{cm} en
\qty{92}{s}. \omterm{def:b1:cell-unit-of-life:prokeuk}{Noyau} à $1000\,g$~: $v = \frac{2}{9}\times(4\times
10^{-6})^2\times 100\times 9810/10^{-3} = \qty{3.5e-3}{m/s}$~;
\qty{5}{cm} en \qty{14}{s}.
\end{solution}

\begin{solution}{exo:b1:cell-unit-of-life:8}
La succinate déshydrogénase est mitochondriale~: \qty{80}{\%} des
mitochondries sont dans le culot 2~; les \qty{15}{\%} du culot 1 sont
des mitochondries piégées dans des \omterm{def:b1:cell-unit-of-life:cell}{cellules} entières et des débris ou
entraînées avec les noyaux~; les \qty{5}{\%} du surnageant proviennent
de mitochondries rompues pendant l'homogénéisation. La lactate
déshydrogénase est cytosolique et soluble~: elle reste dans le
surnageant final.
\end{solution}

\begin{solution}{exo:b1:cell-unit-of-life:9}
(a) Microscopie photonique à contraste de phase sur \omterm{def:b1:cell-unit-of-life:cell}{cellules} vivantes
(le processus dure quelques minutes et exige des \omterm{def:b1:cell-unit-of-life:cell}{cellules} vivantes).
(b) Microscopie photonique sur coupe colorée, ou MET pour être sûr (les
mitochondries sont à la limite du photonique). (c) MET (les pores font
\qty{100}{nm}). (d) MEB (relief de surface).
\end{solution}

\begin{solution}{exo:b1:cell-unit-of-life:10}
Le ballon resté limpide exclut que l'air lui-même, ou le bouillon
bouilli lui-même, engendre la vie~: l'air entre librement par le col
ouvert. Le ballon devenu trouble, après contact avec la poussière,
montre que ce que le col avait retenu --- des particules venues de
l'air --- est la source des \omterm{def:b1:organism-environment:organism}{organismes}. Ensemble, ils prouvent que les
microbes viennent de microbes portés par la poussière. Un ballon
scellé et bouilli laissait ouverte l'objection que le scellement avait
privé le contenu de l'air dont une «~force vitale~» aurait eu besoin~;
le col de cygne admet l'air et n'exclut que la poussière.
\end{solution}

\begin{solution}{exo:b1:cell-unit-of-life:11}
$10^5 = 2^{n}$ donne $n = 16.6$ doublements, soit \qty{415}{min},
environ \qty{7}{h}.
$10^8\times 10^3\,\mathrm{mL}\times\qty{1}{pg} = \qty{0.1}{g}$
par litre. La croissance s'arrête parce que l'oxygène (que le milieu ne
peut fournir qu'à quelques milligrammes par litre), un minéral, ou
l'accumulation d'acides et de déchets devient limitant, et non le
glucose.
\end{solution}

\begin{solution}{exo:b1:cell-unit-of-life:12}
Une \omterm{def:b1:cell-unit-of-life:cell}{cellule} dix fois plus grande dans chaque dimension a dix fois moins
de surface membranaire par unité de volume~; comme les membranes portent
la machinerie de transduction d'énergie et de transport, la grande
\omterm{def:b1:cell-unit-of-life:cell}{cellule} mourrait de faim. L'eucaryote multiplie la membrane en la
repliant à l'intérieur d'elle-même (réticulum endoplasmique, membrane
interne des mitochondries, thylakoïdes) et crée du même coup des
compartiments clos où des processus incompatibles --- digestion acide,
chimie oxydative, conservation de l'ADN --- fonctionnent séparément,
chacun dans ses conditions. L'image est juste pour ce qui touche aux
membranes, et incomplète~: l'eucaryote a aussi acquis un cytosquelette
et un \omterm{def:b1:cell-unit-of-life:prokeuk}{noyau}, et ses mitochondries sont des descendantes de bactéries
englobées, non des replis.
\end{solution}

\begin{solution}{pb:b1:cell-unit-of-life:1}
\textbf{1.} Culot 1~: noyaux (et \omterm{def:b1:cell-unit-of-life:cell}{cellules} entières, débris). Culot 2~:
mitochondries (avec lysosomes et peroxysomes). Culot 3~: microsomes
(fragments de réticulum endoplasmique) et ribosomes.
\textbf{2.} Le froid ralentit les enzymes (dont les protéases et les
phospholipases libérées par la rupture) et préserve les \omterm{def:b1:cell-unit-of-life:prokeuk}{organites}~;
l'isotonie évite que les \omterm{def:b1:cell-unit-of-life:prokeuk}{organites} ne gonflent et n'éclatent, ou ne se
rétractent.
\textbf{3.} Les mitochondries (membrane interne)~; le cytosol.
\textbf{4.} $500 + 200 + 300 + 950 = \qty{1950}{mg}$, soit
\qty{97.5}{\%}~: pertes sur les \omterm{def:b1:cell-unit-of-life:prokeuk}{parois} des tubes et lors des transferts.
\textbf{5.} SDH~: $15 + 60 + 5 + 4 = 84$, soit \qty{84}{\%} (une partie
de l'activité est perdue par lésion des mitochondries). LDH~:
$20 + 12 + 8 + 350 = 390$, soit \qty{97.5}{\%}.
\textbf{6.} $2000/10 = \qty{200}{mg/g}$~: \qty{20}{\%} de la masse fraîche (le reste est surtout de l'eau).
\textbf{7.} $v = \frac{2}{9}\times(5\times 10^{-7})^2\times 100\times
\num{9.81e4}/10^{-3} = \qty{5.4e-4}{m/s}$~; \qty{5}{cm} en \qty{92}{s}~:
\qty{20}{min} sont largement suffisantes.
\textbf{8.} À $1000\,g$, $v = \qty{5.4e-5}{m/s}$~; en \qty{600}{s},
\qty{3.3}{cm}. Réparties uniformément sur \qty{5}{cm}, les mitochondries
situées à moins de \qty{3.3}{cm} du fond l'atteignent~: \qty{65}{\%}, si
rien d'autre n'intervenait.
\textbf{9.} Seuls \qty{15}{\%} sédimentent en fait~: un culot réel ne
retient pas tout ce qui l'atteint (les noyaux tassés sont remis en
suspension avec ménagement et la couche supérieure, lâche, reste dans le
surnageant), et la formule surestime la vitesse d'\omterm{def:b1:cell-unit-of-life:prokeuk}{organites} non
sphériques et hydratés. Des mitochondries incluses dans des \omterm{def:b1:cell-unit-of-life:cell}{cellules}
entières ou adhérant aux noyaux entrent aussi dans le culot 1.
\textbf{10.} $v = \frac{2}{9}\times(4\times 10^{-6})^2\times 100\times
9810/10^{-3} = \qty{3.5e-3}{m/s}$~; \qty{14}{s}.
\textbf{11.} $v = \frac{2}{9}\times(1.2\times 10^{-8})^2\times 600\times
\num{9.81e5}/10^{-3} = \qty{1.9e-5}{m/s}$~; \qty{5}{cm} en
\qty{2650}{s}, soit \qty{44}{min}. À $\num{10000}\,g$ la vitesse vaut
\qty{1.9e-6}{m/s}~: en \qty{20}{min}, \qty{2.3}{mm}, de sorte que les
ribosomes restent dans le surnageant 2.
\textbf{12.} La vitesse varie comme $r^2$~: à faible force, seules les
plus grosses particules se déplacent sensiblement, de sorte que chaque
étape sédimente une classe et laisse les plus petites. Trop courte, la
classe est incomplètement sédimentée~; trop longue, la classe suivante
commence à contaminer le culot.
\textbf{13.} \omterm{def:b1:cell-unit-of-life:fractionation}{Homogénat} $100/2000 = \qty{0.05}{U/mg}$~; culot 1
$15/500 = 0.03$~; culot 2 $60/200 = 0.30$~; culot 3 $5/300 = 0.017$~;
surnageant $4/950 = 0.004$.
\textbf{14.} Purification $0.30/0.05 = 6$~; rendement $60/100 =
\qty{60}{\%}$.
\textbf{15.} $0.017/0.05 = 0.33$~: le culot 3 est appauvri en mitochondries~; le peu d'activité qu'il retient vient de fragments de membrane mitochondriale produits par l'homogénéisation, qui sédimentent avec les microsomes.
\textbf{16.} LDH~: \omterm{def:b1:cell-unit-of-life:fractionation}{homogénat} $400/2000 = \qty{0.2}{U/mg}$~; surnageant
$350/950 = \qty{0.37}{U/mg}$~; purification $1.8$~; rendement \qty{87.5}{\%}.
\textbf{17.} Le cytosol représente la moitié des protéines de la
\omterm{def:b1:cell-unit-of-life:cell}{cellule}, de sorte que même un cytosol pur ne peut que doubler l'activité
spécifique~; les mitochondries en représentent le sixième ou moins, de
sorte que les isoler concentre leurs enzymes six fois ou davantage. Le
rendement est élevé pour l'enzyme soluble parce qu'elle n'est pas perdue
avec les \omterm{def:b1:cell-unit-of-life:prokeuk}{organites} lésés.
\textbf{18.} Non~: le culot 2 retient \qty{200}{mg} de protéines et
piège du surnageant entre ses particules~; \qty{12}{U} font \qty{3}{\%}
du total, ce qui est compatible avec du cytosol piégé. Test~: laver le
culot (remise en suspension et nouvelle centrifugation)~: une enzyme
piégée part, une enzyme liée reste.
\textbf{19.} Si le culot 2 est à \qty{60}{\%} de mitochondries, des
mitochondries pures ont $0.30/0.6 = \qty{0.5}{U/mg}$~; les \qty{100}{U}
du foie correspondent alors à \qty{200}{mg} de protéines
mitochondriales, soit \qty{10}{\%} des \qty{2000}{mg}.
\textbf{20.} MET sur culot fixé, inclus, coupé finement et contrasté aux
métaux lourds~: des profils en haricot à double membrane et à crêtes,
avec quelques lysosomes (corps denses) et peroxysomes parmi eux.
\textbf{21.} Observer une goutte du culot remis en suspension au microscope à contraste de phase et compter les \omterm{def:b1:cell-unit-of-life:cell}{cellules} intactes parmi les noyaux libres~; s'il y en a beaucoup, homogénéiser plus complètement (plus de va-et-vient, piston plus serré) et recentrifuger, car les \omterm{def:b1:cell-unit-of-life:prokeuk}{organites} piégés abaissent tous les rendements.
\textbf{22.} La \omterm{def:b1:cell-unit-of-life:fractionation}{centrifugation différentielle} sépare selon la taille
(vitesse $\propto r^2$) et les trois \omterm{def:b1:cell-unit-of-life:prokeuk}{organites} ont des tailles
voisines~; dans un gradient de densité, une particule s'arrête à sa
propre masse volumique quelle que soit sa taille, et les trois masses
volumiques diffèrent.
\textbf{23.} Des lysosomes intacts gardent l'enzyme derrière une
membrane que le substrat ne franchit pas (latence)~; le détergent rompt
la membrane et l'expose. Les lysosomes du culot sont donc pour la
plupart intacts.
\textbf{24.} L'eau est hypotonique~: les \omterm{def:b1:cell-unit-of-life:prokeuk}{organites} gonflent et éclatent.
Mitochondries et lysosomes se rompent, de sorte que le culot 2 diminue~;
la succinate déshydrogénase (liée aux membranes) passe en partie dans le
culot 3 sous forme de fragments membranaires~; la phosphatase acide
(soluble dans le lysosome) est libérée dans le surnageant final.
\textbf{25.} Facteur de purification $6$, rendement \qty{60}{\%}. Le
facteur est limité par la contamination (lysosomes, peroxysomes, cytosol
piégé~: le culot 2 est à \qty{60}{\%} de mitochondries)~; le rendement,
par les mitochondries perdues dans le culot 1 (piégées) et rompues
pendant l'homogénéisation.
\end{solution}
